\chapter{Introduction}
\label{ch:introduction}

\section{Background}

Deciding when to buy a stock and when to sell it is one of the most common financial decisions
an individual makes, and one of the hardest to make well. The classical view holds that prices
already reflect available information, so that no rule can beat the market after costs
\citep{fama1970efficient}; the adaptive view holds that predictability can appear and fade as
conditions change, so whether an edge exists is an empirical question for each market
\citep{lo2004adaptive}. Machine learning now routinely finds weak but real structure in stock
returns \citep{krauss2017deep,gu2020empirical}. Whether that structure can be turned into
better decisions on the Vietnamese market, under its own rules, is far less studied.

\section{Personal Motivation}

Few people have never dreamed of financial freedom, of money that earns money while its owner
sleeps. I am no exception. The stock market is the most accessible place to chase that dream,
and for a data science student it is also an irresistible puzzle: thousands of new numbers
every day, and the promise that somewhere in them lies a pattern others have missed.

The Vietnamese market adds its own twist. In my experience, know-how here is closely guarded:
those who appear to have an edge rarely explain it, and what is shared in public is mostly tips
without method. Whether real alpha exists on \gls{hose} or not, it is not written down anywhere a
newcomer can learn from. That alone made me want to look for myself.

There is, honestly, some ego involved as well: the wish to find out whether one student with
public data and careful methods can beat a market that many professionals do not. A thesis is a
rare chance to try this seriously, with costs, market rules and statistics strict enough that
the answer is believable whichever way it turns out.

Finally, I want the work to be useful to those who come after me. Most forecasting research is
built on large, developed markets, and its lessons do not carry over cleanly to a market with a
$\pm7\%$ daily band, delayed settlement, opaque speculative operators and data that must be
repaired before use. If this thesis leaves behind a clean dataset, a fair way to evaluate
strategies and an honest record of what works and what does not, the next student who tries to
beat the Vietnamese market can start where this one ends rather than where it began.

\section{Problem Statement}

A trading aid for the \gls{hose} faces weaknesses on three sides. Each one, if ignored, makes a
system look better on paper than it can be in practice.

\paragraph{The market.}
\begin{itemize}
  \item The $\pm7\%$ daily band can lock a falling stock at the floor, so exits are not
    guaranteed; settlement forces a position to be held for several sessions.
  \item Trading concentrates in a small group of stocks. Thin stocks have stale prices and are
    exposed to speculative operators who move prices and paint misleading chart patterns
    \citep{allen1992manipulation,aggarwal2006manipulation}.
  \item Fees on both sides and a tax on every sale erase many small statistical edges.
\end{itemize}

\paragraph{The data.}
\begin{itemize}
  \item Sources built from current listings omit delisted companies (survivorship bias).
  \item Corporate actions are adjusted inconsistently; taken as delivered, a commercial feed
    contains hundreds of price drops that never happened (Section~\ref{sec:data}).
  \item Defects are silent: duplicated rows reported as success, placeholder prices for
    sessions with no trade, fields that disagree with each other.
\end{itemize}

\paragraph{Common approaches.}
\begin{itemize}
  \item Backtests ignore the band, settlement, lot size, fees and taxes, and report the best of
    many configurations without correcting for the number tried
    \citep{harvey2016cross,bailey2014dsr,bailey2017pbo}.
  \item Features and labels leak information not yet known on the day of the decision, and a
    single test period covers only one market regime.
  \item Models output a probability, not a decision: they do not say whether to act, how much,
    when to exit or why, which matters most for investors already prone to overtrading and to
    holding losers \citep{barber2000trading,shefrin1985disposition}.
\end{itemize}

\section{Proposed Direction}

This thesis develops an \gls{ai}-driven \gls{dss} for trading \gls{hose} stocks: a system that
supports the investor's decision rather than replacing it \citep{power2002dss}. In outline:

\begin{itemize}
  \item a cleaned, survivorship-free daily dataset of \gls{hose} stocks, with its defects
    measured and documented;
  \item a signal layer built on machine-learning models and interpretable price--volume
    features;
  \item a decision layer that turns signals into recommendations with reasons, position sizes
    and risk warnings;
  \item an evaluation that enforces \gls{hose} rules and costs, uses walk-forward validation
    and corrects for multiple testing;
  \item an interactive application with historical scenario replays, and an \gls{mt5}
    implementation as an independent check.
\end{itemize}

Research questions, scope and the detailed design are fixed in the following chapters as the
work progresses.

\section{Thesis Structure}

Chapter~\ref{ch:related-work} reviews related work. Chapter~\ref{ch:methodology} describes the
data and the system design. Chapter~\ref{ch:implementation} presents the implementation,
Chapter~\ref{ch:results} the experiments and Chapter~\ref{ch:discussion} their interpretation.
Chapter~\ref{ch:conclusion} concludes.
